\chapter{Lezione 5 --- Trasformata z, esempi elementari e corrispondenza tra piano \(s\) e piano \(z\)}

\section{Richiamo dalla lezione precedente}

Nella lezione precedente abbiamo concluso che, se un segnale continuo causale \(f(t)\)
viene campionato con periodo \(T\), allora il segnale campionato in forma impulsiva può essere scritto come
\[
\hat f(t)=\sum_{k=0}^{+\infty} f(kT)\,\delta(t-kT).
\]

Applicando la trasformata di Laplace si ottiene
\[
\mathcal{L}\{\hat f(t)\}
=
\int_{0^-}^{+\infty}
\left(
\sum_{k=0}^{+\infty} f(kT)\,\delta(t-kT)
\right)e^{-st}\,dt.
\]

Scambiando somma e integrale,
\[
\mathcal{L}\{\hat f(t)\}
=
\sum_{k=0}^{+\infty} f(kT)\int_{0^-}^{+\infty}\delta(t-kT)e^{-st}\,dt
=
\sum_{k=0}^{+\infty} f(kT)e^{-skT}.
\]

Questa espressione è una successione di termini numerici pesati da \(e^{-skT}\),
e suggerisce naturalmente l'introduzione di una nuova variabile complessa.

\section{Nascita della trasformata \(z\)}

Poiché l'esponenziale \(e^{-skT}\) non è particolarmente comodo da manipolare,
si introduce la variabile
\[
z=e^{sT},
\qquad\text{per cui}\qquad
e^{-sT}=z^{-1}.
\]

La trasformata di Laplace del segnale campionato diventa allora
\[
\mathcal{L}\{\hat f(t)\}
=
\sum_{k=0}^{+\infty} f(kT)z^{-k}.
\]

Questa espressione motiva la definizione della \textbf{trasformata \(z\)}.

\section{Definizione della trasformata \(z\)}

\subsection{Trasformata \(z\) monolatera}

Per una sequenza discreta
\[
x[k], \qquad k=0,1,2,\dots
\]
si definisce la trasformata \(z\) monolatera come
\[
X(z)=\mathcal{Z}\{x[k]\}=\sum_{k=0}^{+\infty} x[k]\,z^{-k}.
\]

Nel nostro contesto,
\[
x[k]=f(kT),
\]
quindi
\[
\mathcal{Z}\{f(kT)\}=\sum_{k=0}^{+\infty} f(kT)\,z^{-k}.
\]

La trasformata \(z\) monolatera è l'analogo discreto della trasformata di Laplace unilaterale:
parte da \(k=0\) e arriva a \(+\infty\), quindi è adatta a segnali causali e transitori.

\subsection{Trasformata \(z\) bilatera}

Per sequenze definite su tutti gli interi, si introduce invece la trasformata \(z\) bilatera:
\[
X(z)=\mathcal{Z}_b\{x[k]\}=\sum_{k=-\infty}^{+\infty} x[k]\,z^{-k}.
\]

Questa forma è l'analogo discreto della trasformata di Fourier/Laplace bilatera
ed è utile per segnali stazionari o definiti su tutto l'asse discreto.

\subsection{Osservazione}

La sequenza \(x[k]\) è discreta, ma \(X(z)\) è una funzione della variabile complessa
\[
z\in\mathbb{C}.
\]
Quindi la trasformata \(z\) non è un oggetto discreto:
\[
X(z)
\]
è una funzione continua della variabile complessa \(z\), definita sul proprio dominio di convergenza.

\section{Interpretazione di \(z^{-1}\) e \(z\)}

Negli appunti compare una lettura operativa molto utile:
\begin{itemize}
    \item \(z^{-1}\) si interpreta come \textbf{operatore di ritardo unitario};
    \item \(z\) si interpreta formalmente come \textbf{operatore di anticipo unitario}.
\end{itemize}

Questa interpretazione è fondamentale nei sistemi discreti e nei diagrammi a blocchi.

\section{Sequenze equispaziate}

La trasformata \(z\) vale per sequenze discrete i cui campioni siano equispaziati nel tempo,
cioè prelevati con periodo costante \(T\).

Dire che i campioni sono equispaziati significa che tra due campioni consecutivi
passa sempre lo stesso intervallo di tempo:
\[
t_k = kT.
\]

Questo è il caso naturale che si incontra nel campionamento uniforme.

\section{Proprietà della trasformata \(z\)}

La trasformata \(z\) monolatera gode di proprietà del tutto analoghe a quelle della Laplace.

\subsection{Linearità}

Per \(\alpha,\beta\in\mathbb{R}\) oppure \(\mathbb{C}\),
\[
\mathcal{Z}\{\alpha x[k]+\beta y[k]\}
=
\alpha X(z)+\beta Y(z).
\]

\subsection{Convoluzione discreta}

Se \(x[k]\) e \(y[k]\) sono due sequenze causali, la loro convoluzione discreta è
\[
(x\ast y)[n]
=
\sum_{k=0}^{n} x[k]\,y[n-k].
\]

La trasformata \(z\) trasforma la convoluzione in prodotto:
\[
\mathcal{Z}\{x\ast y\}=X(z)Y(z).
\]

Questa proprietà è perfettamente analoga a quella già vista per Laplace nel caso continuo.

\section{Primo esempio: trasformata \(z\) del gradino}

Consideriamo il gradino continuo
\[
f(t)=1(t)=H(t),
\]
e adottiamo la convenzione coerente con gli appunti secondo cui il segnale campionato vale
\[
f(kT)=1,
\qquad k=0,1,2,\dots
\]

Allora
\[
X(z)=\sum_{k=0}^{+\infty} 1\cdot z^{-k}
=
\sum_{k=0}^{+\infty}(z^{-1})^k.
\]

Questa è una serie geometrica, che converge se
\[
|z^{-1}|<1
\qquad\Longleftrightarrow\qquad
|z|>1.
\]

Quindi
\[
X(z)=\frac{1}{1-z^{-1}}
=
\frac{z}{z-1},
\qquad
\text{ROC: } |z|>1.
\]

\subsection{Osservazioni}

\begin{itemize}
    \item La trasformata è una funzione razionale.
    \item C'è un polo in \(z=1\).
    \item C'è uno zero in \(z=0\), se si guarda la forma \(\dfrac{z}{z-1}\).
\end{itemize}

Questo mostra che, anche nel dominio \(z\), possiamo continuare a ragionare in termini di poli,
zeri e funzioni razionali.

\section{Regione di convergenza}

Come nella trasformata di Laplace, anche la trasformata \(z\) non è caratterizzata solo
dalla sua espressione algebrica, ma anche dalla \textbf{regione di convergenza} (ROC).

Nel primo esempio abbiamo ottenuto
\[
X(z)=\frac{z}{z-1},
\qquad \text{ROC: } |z|>1.
\]

Questo significa che la funzione è definita e utile nella regione esterna al cerchio di raggio \(1\),
escluso il polo \(z=1\).

\subsection{Osservazione sul dominio}

Negli appunti compare l'idea, mutuata dall'analisi matematica, di studiare i punti di bordo
del dominio guardando i limiti dai punti in cui la funzione è effettivamente definita.

Questa osservazione serve a ricordare che:
\begin{itemize}
    \item una funzione razionale può essere prolungata per continuità in un punto
    solo se il limite esiste ed è finito;
    \item in presenza di un polo il prolungamento non è possibile;
    \item quindi il punto resta escluso dal dominio.
\end{itemize}

\section{Secondo esempio: trasformata \(z\) di un esponenziale campionato}

Consideriamo il segnale continuo causale
\[
f(t)=e^{-at}1(t),
\qquad a>0.
\]

Campionandolo agli istanti \(t=kT\), otteniamo
\[
f(kT)=e^{-akT},
\qquad k=0,1,2,\dots
\]

La trasformata \(z\) è allora
\[
X(z)=\sum_{k=0}^{+\infty} e^{-akT}z^{-k}
=
\sum_{k=0}^{+\infty}\left(e^{-aT}z^{-1}\right)^k.
\]

La serie converge se
\[
|e^{-aT}z^{-1}|<1
\qquad\Longleftrightarrow\qquad
|z|>e^{-aT}.
\]

Quindi
\[
X(z)=\frac{1}{1-e^{-aT}z^{-1}}
=
\frac{z}{z-e^{-aT}},
\qquad
\text{ROC: } |z|>e^{-aT}.
\]

\subsection{Interpretazione}

La trasformata ha:
\begin{itemize}
    \item un polo in
    \[
    z=e^{-aT},
    \]
    \item uno zero in \(z=0\), nella forma razionale in \(z\).
\end{itemize}

Poiché \(a>0\), si ha
\[
0<e^{-aT}<1,
\]
quindi il polo cade all'interno del cerchio unitario.

\section{Funzioni razionali nel dominio \(z\)}

Un punto sottolineato negli appunti è che la trasformata \(z\) dei segnali elementari
porta ancora a funzioni razionali.

Questo è importante perché permette di mantenere, anche nel discreto,
molti degli strumenti concettuali usati nel continuo:
\begin{itemize}
    \item poli;
    \item zeri;
    \item fattorizzazioni;
    \item studio della stabilità;
    \item diagrammi a blocchi.
\end{itemize}

\subsection{Osservazione tecnica}

Nel dominio \(z\) occorre sempre essere attenti alla convenzione usata:
una funzione può essere scritta come polinomio in \(z\) oppure come polinomio in \(z^{-1}\).
Per questo alcune considerazioni su grado del numeratore e del denominatore
vanno interpretate con cautela, specificando sempre quale forma si sta usando.

\section{Corrispondenza tra piano \(s\) e piano \(z\)}

La relazione fondamentale tra piano continuo e piano discreto è
\[
z=e^{sT}.
\]

Se
\[
s=\sigma+j\omega,
\]
allora
\[
z=e^{(\sigma+j\omega)T}=e^{\sigma T}e^{j\omega T}.
\]

Questa formula contiene tutta la corrispondenza geometrica tra i due piani.

\subsection{Caso \(\omega=0\)}

Se \(\omega=0\), allora
\[
s=\sigma
\qquad\Longrightarrow\qquad
z=e^{\sigma T}.
\]

Quindi:
\begin{itemize}
    \item se \(\sigma=0\), allora \(z=1\);
    \item se \(\sigma>0\), allora \(z=e^{\sigma T}>1\), cioè si va sull'asse reale positivo
    all'esterno del cerchio unitario;
    \item se \(\sigma<0\), allora \(0<z=e^{\sigma T}<1\), cioè si va sull'asse reale positivo
    all'interno del cerchio unitario;
    \item se \(\sigma\to -\infty\), allora \(z\to 0\).
\end{itemize}

\subsection{Caso \(\sigma=0\)}

Se \(\sigma=0\), allora
\[
s=j\omega
\qquad\Longrightarrow\qquad
z=e^{j\omega T}=\cos(\omega T)+j\sin(\omega T).
\]

Quindi l'asse immaginario del piano \(s\) viene mappato sul cerchio unitario nel piano \(z\):
\[
|z|=1.
\]

\subsection{Osservazione importante}

Questa mappa non è biunivoca su tutto l'asse immaginario:
infatti
\[
e^{j\omega T}=e^{j(\omega+2k\pi/T)T},
\qquad k\in\mathbb{Z}.
\]

Cioè lo stesso punto del piano \(z\) corrisponde a infiniti valori di \(\omega\)
che differiscono di multipli di \(2\pi/T\).

\section{Caso generale}

Nel caso generale
\[
z=e^{\sigma T}e^{j\omega T}.
\]

Da qui si leggono due proprietà geometriche fondamentali:

\begin{itemize}
    \item fissato \(\sigma\), il modulo è
    \[
    |z|=e^{\sigma T},
    \]
    quindi le rette verticali \(\sigma=\text{costante}\) nel piano \(s\)
    si trasformano in circonferenze concentriche nel piano \(z\);

    \item fissato \(\omega\), l'argomento è
    \[
    \arg(z)=\omega T \pmod{2\pi},
    \]
    quindi le rette orizzontali \(\omega=\text{costante}\) nel piano \(s\)
    si trasformano in semirette uscenti dall'origine nel piano \(z\).
\end{itemize}

\section{Conseguenze fondamentali della mappa \(z=e^{sT}\)}

Dalla formula precedente segue immediatamente che:

\begin{itemize}
    \item l'asse immaginario del piano \(s\) diventa il cerchio unitario del piano \(z\);
    \item il semipiano sinistro
    \[
    \Re(s)<0
    \]
    viene mappato all'interno del cerchio unitario
    \[
    |z|<1;
    \]
    \item il semipiano destro
    \[
    \Re(s)>0
    \]
    viene mappato all'esterno del cerchio unitario
    \[
    |z|>1;
    \]
    \item l'origine del piano \(s\), cioè \(s=0\), viene mappata nel punto
    \[
    z=1.
    \]
\end{itemize}

\section{Fascia primaria}

Per rendere biunivoca la corrispondenza tra \(s\) e \(z\), bisogna limitarsi a una fascia
del piano \(s\) di ampiezza \(2\pi/T\) rispetto alla variabile \(\omega\).

Una scelta standard è la \textbf{fascia primaria}
\[
-\frac{\pi}{T}<\omega\le \frac{\pi}{T}.
\]

Se ci si limita a questa fascia, la mappa
\[
z=e^{sT}
\]
diventa univoca.

\subsection{Interpretazione}

La massima pulsazione contenuta nella fascia primaria è
\[
\omega_{\max}=\frac{\pi}{T}.
\]

Se definiamo la pulsazione di campionamento
\[
\omega_c=\omega_s=\frac{2\pi}{T},
\]
allora
\[
\omega_{\max}=\frac{\omega_c}{2}.
\]

Questa è un'altra forma del teorema del campionamento:
la fascia primaria contiene le frequenze non aliasate.

\section{Legame con il teorema di Shannon}

Poiché la fascia primaria è
\[
-\frac{\pi}{T}<\omega\le \frac{\pi}{T},
\]
per evitare aliasing bisogna fare in modo che il contenuto in frequenza del segnale
ricada interamente dentro questa fascia.

Equivalentemente:
\[
\omega_b < \frac{\pi}{T}
\qquad \Longleftrightarrow \qquad
\omega_s=\frac{2\pi}{T}>2\omega_b.
\]

Questa è esattamente la formulazione del teorema di Shannon nel dominio delle pulsazioni.

\section{Mappatura di poli e singolarità}

Un altro punto importante richiamato negli appunti è che le singolarità della trasformata
nel dominio continuo vengono mappate nei corrispondenti punti del piano \(z\).

Se un sistema continuo ha un polo in
\[
s=p,
\]
allora il corrispondente punto nel dominio discreto è
\[
z=e^{pT}.
\]

Questo significa che:
\begin{itemize}
    \item poli continui nel semipiano sinistro vanno all'interno del cerchio unitario;
    \item poli continui sull'asse immaginario vanno sul cerchio unitario;
    \item poli continui nel semipiano destro vanno all'esterno del cerchio unitario.
\end{itemize}

Questa corrispondenza è alla base del criterio di stabilità per i sistemi discreti.

\section{Sintesi della lezione}

I punti principali della lezione sono:
\begin{enumerate}
    \item derivazione della trasformata \(z\) dalla trasformata di Laplace del segnale campionato;
    \item definizione di trasformata \(z\) monolatera e bilatera;
    \item interpretazione di \(z^{-1}\) come operatore di ritardo;
    \item proprietà di linearità e convoluzione;
    \item calcolo della trasformata \(z\) del gradino:
    \[
    \mathcal{Z}\{1\}=\frac{z}{z-1},\qquad |z|>1;
    \]
    \item calcolo della trasformata \(z\) dell'esponenziale campionato:
    \[
    \mathcal{Z}\{e^{-akT}\}=\frac{z}{z-e^{-aT}},\qquad |z|>e^{-aT};
    \]
    \item corrispondenza tra piano \(s\) e piano \(z\) tramite
    \[
    z=e^{sT};
    \]
    \item asse immaginario di \(s\) mappato sul cerchio unitario;
    \item semipiano sinistro mappato all'interno del cerchio unitario;
    \item introduzione della fascia primaria
    \[
    -\frac{\pi}{T}<\omega\le \frac{\pi}{T},
    \]
    che rappresenta la porzione non aliasata del piano \(s\).
\end{enumerate}